\documentclass[border=10pt,tikz]{standalone}
\usepackage{tikz}
\usepackage{tikz-timing}
\usepackage{fontspec}
\usepackage{xcolor}
\setmainfont{Apple SD Gothic Neo}
\setsansfont{Apple SD Gothic Neo}
\definecolor{warnfill}{HTML}{FECACA}
\definecolor{okfill}{HTML}{BBF7D0}
\definecolor{hifill}{HTML}{FED7AA}
\usetikztiminglibrary{counters,interval,nicetabs}
\begin{document}
\begin{tikzpicture}[font=\scriptsize]

\node[font=\footnotesize\bfseries, anchor=west] at (0, 3) {I²C Transaction — Start + 7-bit Address + Write + ACK + Stop};

\node[anchor=west] at (0, 0) {\begin{tikztimingtable}
  SCL  & 0.5H 0.5L 0.5H 0.5L 0.5H 0.5L 0.5H 0.5L 0.5H 0.5L 0.5H 0.5L 0.5H 0.5L 0.5H 0.5L 0.5H 0.5L 0.5H 0.5L 0.5H \\
  SDA  & 0.5H 0.25L 0.25L 0.5L 0.5H 0.5L 0.5L 0.5H 0.5H 0.5L 0.5L 0.5H 0.5L 0.5H 0.5L 0.5L 0.5L 0.5L 0.5L 0.5H 0.5H \\
\end{tikztimingtable}};

\node[anchor=west, font=\tiny, align=left] at (0, -4) {
  \textbf{S} (Start) = SCL High에서 SDA Hi→Lo\\
  Addr = 7 bit, 다음 1 bit = R/W (이 예에서 W = 0)\\
  \textbf{ACK} = 9번째 SCL pulse 시 슬레이브가 SDA Low\\
  \textbf{P} (Stop) = SCL High에서 SDA Lo→Hi
};

\end{tikzpicture}
\end{document}
